\documentclass[10pt,a4paper]{amsart}
\usepackage[margin=19mm]{geometry}
\usepackage{amsmath,amssymb,mathtools}
\usepackage[scr=rsfs,cal=euler]{mathalfa}
\usepackage{xfrac}
\usepackage[unicode,hidelinks,bookmarks=false]{hyperref}
\newcommand{\ie}{i.e.\ }
\newcommand{\cf}{cf.\ }
\newcommand{\resp}{respectively\ }
\newcommand{\Subsection}{\subsection}
\providecommand{\nobreakdash}{\nobreak\hbox{-}}
\setlength{\emergencystretch}{2em}
\multlinegap0pt
\usepackage{bm}
\usepackage{bbm}
\usepackage{dsfont}
\usepackage{xspace}
\usepackage{cancel}
\usepackage{mathtools}
\usepackage{amsthm}
\makeatletter
\@ifundefined{theorem}{\newtheorem{theorem}{Theorem}}{}
\@ifundefined{corollary}{\newtheorem{corollary}[theorem]{Corollary}}{}
\makeatother
\title[On the global breadth of finite groups with nontrivial parti]{On the global breadth of finite groups with nontrivial partitions}
\author{Seid Kassaw Muhie, Daniele Ettore Otera, Francesco G.Russo}
\date{Source: arXiv:2506.19675v1 (CC BY 4.0)}
\begin{document}
\maketitle

\noindent\emph{Source and licence.} On the global breadth of finite groups with nontrivial partitions. Version arXiv:2506.19675v1 (CC BY 4.0), distributed under the Creative Commons Attribution 4.0 International licence, as recorded in the arXiv licence metadata for this exact version and shown on the abstract page https://arxiv.org/abs/2506.19675v1. Full rights notice: licenses/arxiv-2506.19675-CC-BY-4.0.txt.\par\smallskip

\noindent\textit{Editorial scope.} Counted unit: the Corollary whose source label is \texttt{l2} in the article (source0.tex lines 558--561, source label \texttt{l2}), delivered together with the article's own complete original proof of it (source0.tex lines 563--578, one \texttt{proof} environment of 16 lines). No mathematics is written, repaired or re-derived: statement and proof are the article's own text line for line. Required context retained uncounted: l0 (lines 181--182); l1 (lines 484--490). Every definition, display and label of those lines is retained unchanged. Omitted from this package and not redistributed: 1--180, 183--483, 491--557, 562--562, 579--991. Nothing inside a retained excerpt is omitted or altered: every retained body is delivered exactly as the article's lines are written, so no omission receipt of any kind accompanies this package. Citations: the retained excerpts carry no citation command at all, so no bibliography entry of the article is transcribed and no reference is redistributed. Reference adaptation: none was needed and none was made; every reference inside the delivered unit resolves inside it, so no unresolved reference or citation is produced. Printed slips kept literal: no printed word, symbol, hypothesis or number of the counted statement or of its proof was altered. Layout and class: the article's own class is \texttt{amsart} with packages including vmargin, amsmath, amsfonts, amssymb, amsthm, xcolor, amscd, latexsym, xspace, enumerate, hyperref, pdfpages, tikz-cd, pgfplots; the wrapper is the lane's pinned \texttt{amsart} editorial wrapper at 10pt on A4, which restates the theorem-like environments the retained text needs and adds the article's own macro definitions that the retained text needs (none), with extra wrapper packages (bm, bbm, dsfont, xspace, cancel, mathtools). The delivered numbering is the unit's own. Assets: no retained span contains a figure environment, a graphics-inclusion command, a PDF-page inclusion, a table or a TikZ picture; no image file is copied into this package, the package declares no asset dependency and no image file is redistributed with it. Licence: version arXiv:2506.19675v1 of this article is distributed under the Creative Commons Attribution 4.0 International licence, as recorded in the arXiv licence metadata for this exact version and shown on the abstract page https://arxiv.org/abs/2506.19675v1. A full rights notice is delivered at licenses/arxiv-2506.19675-CC-BY-4.0.txt. Characters: every retained excerpt is pure ASCII, so the delivered engine, XeLaTeX with its default Latin Modern text font, reports no missing character.
\begin{theorem}\label{l0} If $G$ is a group with nontrivial partition, then there exists a maximal cyclic subgroup $M$ of $G$ such that $\mathbf{B}(G)=\mathbf{b}_{|M|}(G).$
\end{theorem} 

\begin{corollary}\label{l1}For $q=3^n$ with $n \ge 2$ or for all odd prime powers $q \ge 5$
 \begin{equation} \label{pslodd}
 \mathbf{B}(\mathrm{PSL}(2,q)) + 2 = \frac{q(q-1)}{2},
 \end{equation} but in case $q=2^m$ (with $m \ge 2$)
 \begin{equation}\mathbf{B}(\mathrm{PSL}(2,q)) = 2^{2m-1}.
 \end{equation}
\end{corollary}

\begin{corollary}\label{l2} Let $q=2^{2n+1}$ for $n \ge 1$. Then,
\begin{equation}\mathbf{B}(\mathrm{Sz}(2^{2n+1}))= \frac{q^4-q^3-2(q+1)}{2}.
 \end{equation} 
\end{corollary}

\begin{proof} Arguing as in the proof of Theorem \ref{l0}, we consider $q=2^{2n+1}$ for $n \ge 0$ and know that $|\mathrm{Sz}(q)|= q^2(q^2+1)(q-1)$. Counting the elements of order $q-1$ in $\mathrm{Sz}(q)$, we get
\begin{equation}
|L_{q-1}(\mathrm{Sz}(q))|=|\{x \in \mathrm{Sz}(q) \ | \ x^{q-1}=1\}|=(q-1) \ \ \left(\frac{q^4-q^3-2(q-1)}{2}\right)
\end{equation}
and so \begin{equation}\mathbf{b}_{q-1}(\mathrm{Sz}(q))=\frac{q^4-q^3-2(q+1)}{2}.
\end{equation}
Of course if $n=0$, then $\mathrm{Sz}(2)$ is a group of order 20, which is nonsimple. Assume $q=2^{2n+1}$ for $n \ge 1$ and we should check that among $\mathbf{b}_c(\mathrm{Sz}(q))$ with $c\in \mathrm{Div}(\mathrm{Sz}(q))$ the value $\mathbf{b}_{q-1}(\mathrm{Sz}(q))$ indeed realizes the global breadth $\mathbf{B}(\mathrm{Sz}(q))$. We follow the same procedure of the proof of Corollary \ref{l1}, omitting the details. We just report the following computations, which are useful to identify the main local breadths that one needs. In fact we  only need to check the additional values $\mathbf{b}_{c}(\mathrm{Sz}(q))$ when $c=4$, $c=q+2^{n+1} + 1$ and $c=q-2^{n+1} + 1$, that is, 
\begin{equation}
    \mathbf{b}_4(\mathrm{Sz}(q))=  2^{6n+4}, \ \ \ \ \ \ 
    \mathbf{b}_{q-2^{n+1}+1}(\mathrm{Sz}(q))=  2^{6n+4}+2^{4n+3}+2^{2n+1}+2^{n+1}+1,
\end{equation}  
\begin{equation}
\mathbf{b}_{q+2^{n+1} + 1}(\mathrm{Sz}(q))= 2^{6n+3}+2^{2n+1}-2^{n+1}+1.
\end{equation}
Therefore we get what we claimed for $q=2^{2n+1}$  with $n \geq 1$, namely \begin{equation}\mathbf{B}(\mathrm{Sz}(q)) = \frac{q^4-q^3-2(q+1)}{2}=\:\left(2^{2n+1}\right)^4-\left(2^{2n+1}\right)^3-2\left(2^{2n+1}+1\right).\end{equation} 
\end{proof}

\end{document}
